\section{The wall actions in integral coordinates}\label{sec:actions}

The compact geometric actions were proved in
Proposition~\ref{prop:marked-wall}. We now give their
coordinates on every degree of $R$ and on the normalized
indices. This also verifies compatibility with the
attachment formulas.

\subsection{All ring generators}
At a divisorial wall the action in $R$ is the conjugate
of $s_{a_1}$ by $g$. Put
\begin{equation}\label{eq:reflection-root}
 k_0=2H-D_2,\qquad r=\mu(k_0)+\delta
             =\theta_{v'}(ga_1),\qquad q(k_0)=0,
 \quad q(r)=-2.
\end{equation}
Write $\mathsf S$ for this integral ring automorphism. Its degree-two
action is
\begin{equation}\label{eq:S2}
 \mathsf S(\mu(b))=\mu(b)+q(b,k_0)r,\qquad \mathsf S(\delta)=\delta-2r.
\end{equation}
Set $\mu_i'=\mathsf S(\mu_i)$, $\delta'=\mathsf S(\delta)$. The full
degree-four action is
\begin{equation}\label{eq:S4}
 \begin{aligned}
 A'&=\frac{c_2(Y)+3(\delta')^2}{24}
     =A+\frac{\mu(k_0)^2+\delta\mu(k_0)}2,\\
 T_i'&=\mu_i'\delta',\qquad
 C_{ij}'=\mu_i'\mu_j'-q_{ij}A',\\
 Q_i'&=\frac{(\mu_i')^2-q_{ii}A'-T_i'}2.
 \end{aligned}
\end{equation}
In degree six the images are
$\mathsf S(R_i)=\mathsf S(\mu(\alpha^i))A'$ and
$\mathsf S(R_\delta)=-\delta'A'$. The unit and oriented point are
fixed. Formula \eqref{eq:S4} follows from preservation of
$c_2(Y)$ and \eqref{eq:degree2-products}. Its apparently
rational expressions are integral because they are images
under the integral ring map already constructed; the
second formula for $A'$ also shows its integrality
directly. These identities give a finite coordinate
calculation on the entire lattice, including the classes
outside its integral symmetric square.

They imply
\begin{equation}\label{eq:planes-reflection}
 \mathsf S(p_0)=p_2,\quad \mathsf S(p_2)=p_0,\quad \mathsf S(c_2(Y))=c_2(Y),
 \quad \mathsf S^2=\id.
\end{equation}
Consequently $\mathsf S$ preserves the exceptional ideal. In the
fixed two-plane torsion convention,
\[
 \kappa(\mathsf S z)=5\kappa(z)\pmod8,
\]
because $p_2-3p_0\equiv5(p_0-3p_2)$ modulo eight in
the integral lattice. The ruled conventions in
\eqref{eq:ruled-rho} incorporate exactly this change.
At a plane wall every generator has the identity action
in $R$, with the geometric surface change
\eqref{eq:surface-transform} still understood. It is the
exceptional product in \eqref{eq:flop-correspondence}
that makes this the action of a multiplicative
correspondence.

\subsection{The fixed determinant normalizer}
Let $h=\id$ or $s_{a_1}$ act between the source and target
Mukai lattices, and retain the original $y$ on both sides.
Markman's universal-class identity, followed by the
rank-one normalization argument, gives
\begin{equation}\label{eq:wall-indices}
 \gamma_h\bigl(c_k(I_z^y)\bigr)=
 c_k\left(I_{hz}^y\otimes(\det I_{hy}^y)^{-1}\right).
\end{equation}
Indeed transport first gives the target universal family
normalized by $hy$ rather than by $y$. Starting from the
$y$-normalized target family, its $hy$-index still has
rank one, so tensoring by its inverse determinant is
exactly that change of normalization. Equality of
characters implies equality of integral Chern classes
by torsion-freeness, as in
Proposition~\ref{prop:indices}.

For $h=\id$ the determinant is trivial. For $h=s_{a_1}$,
$(y,a_1)=-1$ and hence $hy=y-a_1$. Additivity of the
index and $\det I_y^y=1$ give
\begin{equation}\label{eq:normalizer-line}
 c_1\left((\det I_{y-a_1}^y)^{-1}\right)=\theta_v(a_1),
 \qquad \Phi\theta_v(a_1)=r.
\end{equation}
Thus in the common ring the transformed character is
$\ch(I_{s_{a_1}z}^y)\exp(r)$, not merely the character
with its input vector reflected. Formula
\eqref{eq:tautological-character} and the Newton identities
calculate all its Chern classes. Functoriality includes
these normalizing lines, so they introduce no ambiguity
in the six-step cohomological identity.

\subsection{Dependence on the choices}
An integral change of basis of $\Lambda$ changes the
coordinates in \eqref{eq:basis4} and \eqref{eq:basis6}
by the corresponding integral transformation. The
diagonal construction itself is intrinsic. The isometry
$g$ and the universal normalizer $y$ are specified:
Markman's functoriality fixes the transport, and
\eqref{eq:wall-indices} describes any change of the
normalizer. Different geometric realizations of an
action with the same degree-two marking have the same
cohomological action, by rational generation and
torsion-freeness.

Isotopic regular neighborhoods give the same restriction
and relative maps by naturality of excision. The outward
orientation and the ordered difference $r_V-r_W$ fix
the connecting signs. Replacing a torsion generator by
a unit multiple changes the printed mod-eight functional
by that unit. These choices specify the calculated
cohomology and its maps; no assertion about a homotopy
classification of the gluing follows from them.

\begin{proof}[Proof of Theorem~\ref{thm:main}]
Proposition~\ref{prop:compact-walls} constructs the
compact partial-wall maps, and
Lemma~\ref{lem:opposite-divisors} and
Proposition~\ref{prop:marked-wall} identify their
integral actions. Equations \eqref{eq:Phi}--\eqref{eq:gc}
give the specified ambient marking and exceptional
coordinates. Propositions~\ref{prop:complement},
\ref{prop:boundary} and \ref{prop:all-attachments}
calculate the complement and every attachment map.
Proposition~\ref{prop:indices} and
\eqref{eq:wall-indices} give the universal-index
markings and their determinant correction. Finally
Corollary~\ref{cor:gallery} proves the cohomological
path relation. All computations use the same compact
moduli problem and the same primitive cycles.
\end{proof}

\subsection{Exact verification}
The formulas above are finite expressions in the marked
K3 Gram matrix. Accompanying standard-library Python
programs implement \eqref{eq:blowup} with rational
arithmetic. They compute the integral middle pairing,
the two plane coordinates, the four Chern classes of
each index for a $24$-element topological $K$-basis, and
the images of every degree-two, degree-four and
degree-six basis element under $\mathsf S$. Additional virtual
classes test negative ranks and the determinant
normalization. The retained outputs check the plane
binomial restrictions, all six degree-two attachment
products, and the complete wall loop. The printed
formulas define the calculation independently of these
programs; the compact wall and topological-kernel
arguments are the geometric justifications for using
them.
